\documentclass[tikz,border=3mm,multi=tikzpicture]{standalone}
\usepackage{eleve-fisica}

\begin{document}

% FIGURA 1 — fig-01-interacao-entre-dois-corpos
\begin{tikzpicture}[fisica figura, x=1cm, y=1cm]
  \path[use as bounding box] (-0.2,-0.2) rectangle (9.2,5.8);
  \draw[fisica superficie] (0.5,0.75) -- (8.7,0.75);
  \node[fisica corpo, minimum width=1.9cm, minimum height=1.25cm] (mao)
    at (1.75,1.45) {mão};
  \node[fisica corpo, minimum width=2.25cm, minimum height=1.75cm] (caixa)
    at (7.25,1.65) {caixa};
  \draw[fisica forca] (mao.east) -- (caixa.west)
    node[fisica rotulo, midway, above=8pt] {$\vec F_{\text{mão}\to\text{caixa}}$};
  \node[font=\Large\bfseries, text=eleveazul] at (4.5,4.65)
    {interação};
  \node[font=\large\bfseries, text=elevecinza] at (1.75,3.55)
    {quem age};
  \node[font=\large\bfseries, text=elevecinza] at (7.25,3.55)
    {quem recebe};
  \draw[fisica auxiliar, ->] (1.75,3.2) -- (1.75,2.35);
  \draw[fisica auxiliar, ->] (7.25,3.2) -- (7.25,2.65);
\end{tikzpicture}

% FIGURA 2 — fig-02-intensidade-direcao-e-sentido
\begin{tikzpicture}[fisica figura, x=1cm, y=1cm]
  \path[use as bounding box] (-0.2,-0.2) rectangle (10.1,6.4);
  \coordinate (A) at (1.2,2.75);
  \coordinate (B) at (8.85,2.75);
  \draw[fisica auxiliar] (0.65,2.75) -- (9.45,2.75);
  \draw[fisica forca] (A) -- (B)
    node[fisica rotulo, midway, above=12pt] {$\vec F=20\,\mathrm{N}$};
  \fill[eleveazul] (A) circle (3pt);
  \node[font=\large\bfseries, text=elevecinza, align=center] at (4.95,5.35)
    {intensidade\\comprimento e valor};
  \draw[fisica auxiliar, ->] (4.95,4.65) -- (4.95,3.35);
  \node[font=\large\bfseries, text=elevecinza] at (2.15,0.65)
    {direção horizontal};
  \draw[fisica auxiliar, ->] (3.95,0.85) -- (5.55,2.2);
  \node[font=\large\bfseries, text=elevecinza] at (8.0,0.65)
    {sentido};
  \draw[fisica movimento] (7.25,0.95) -- (8.75,0.95)
    node[fisica rotulo, right] {direita};
\end{tikzpicture}

% FIGURA 3 — fig-03-forcas-de-contato-e-a-distancia
\begin{tikzpicture}[fisica figura, x=1cm, y=1cm]
  \path[use as bounding box] (-0.2,-0.2) rectangle (9.6,10.3);
  \node[font=\Large\bfseries, text=eleveazul] at (4.7,9.75) {contato};
  \draw[fisica superficie] (0.8,6.35) -- (8.7,6.35);
  \node[fisica corpo, minimum width=1.55cm, minimum height=1.1cm] (mao)
    at (1.8,7.0) {mão};
  \node[fisica corpo, minimum width=1.9cm, minimum height=1.45cm] (c1)
    at (6.7,7.1) {caixa};
  \draw[fisica forca] (mao.east) -- (c1.west)
    node[fisica rotulo, midway, above=7pt] {$\vec F$};
  \node[font=\large\bfseries, text=elevecinza] at (4.7,5.55)
    {os corpos se tocam};

  \node[font=\Large\bfseries, text=eleveazul] at (4.7,4.45) {à distância};
  \node[fisica corpo, minimum width=1.55cm, minimum height=2.0cm] (ima)
    at (1.8,2.1) {ímã};
  \node[fisica corpo, circle, minimum size=0.9cm] (clipe)
    at (7.3,2.1) {};
  \draw[eleveazul, line width=1.4pt] (7.05,1.75) -- (7.55,2.45);
  \draw[eleveazul, line width=1.4pt] (7.05,2.45) -- (7.55,1.75);
  \draw[fisica forca] (clipe.west) -- ++(-2.75,0)
    node[fisica rotulo, midway, above=7pt] {$\vec F_m$};
  \node[font=\large\bfseries, text=elevecinza] at (4.7,0.6)
    {há espaço entre os corpos};
\end{tikzpicture}

% FIGURA 4 — fig-04-peso-aponta-para-o-centro
\begin{tikzpicture}[fisica figura, x=1cm, y=1cm]
  \path[use as bounding box] (-0.25,-0.25) rectangle (9.75,9.75);
  \coordinate (O) at (4.75,4.75);
  \shade[ball color=eleveazulclaro] (O) circle (2.75);
  \draw[eleveazul, line width=1.4pt] (O) circle (2.75);
  \fill[eleveazul] (O) circle (3pt);
  \node[fisica rotulo, below=5pt] at (O) {centro};
  \foreach \ang/\pos in {90/above,0/right,270/below,180/left}{
    \coordinate (C) at ($(O)+(\ang:3.45)$);
    \node[fisica corpo, circle, minimum size=0.62cm] at (C) {};
    \draw[fisica forca] (C) -- ($(O)+(\ang:0.95)$);
  }
  \node[fisica rotulo] at (6.25,7.65) {$\vec P$};
  \node[fisica rotulo] at (7.95,3.15) {$\vec P$};
  \node[fisica rotulo] at (3.15,1.7) {$\vec P$};
  \node[fisica rotulo] at (1.55,6.25) {$\vec P$};
\end{tikzpicture}

% FIGURA 5 — fig-05-atrito-se-opoe-ao-deslizamento
\begin{tikzpicture}[fisica figura, x=1cm, y=1cm]
  \path[use as bounding box] (-0.2,-0.2) rectangle (10.2,9.4);
  \node[font=\Large\bfseries, text=eleveazul] at (5.0,8.85)
    {deslizamento para a direita};
  \draw[fisica superficie] (0.65,5.75) -- (9.45,5.75);
  \node[fisica corpo, minimum width=2.0cm, minimum height=1.35cm] (A)
    at (5.0,6.45) {};
  \draw[fisica movimento] ($(A.north)+(0,0.55)$) -- ++(2.65,0)
    node[fisica rotulo, right] {$\vec v$};
  \draw[fisica forca] (A.center) -- ++(-2.65,0)
    node[fisica rotulo, left] {$\vec F_{at}$};

  \node[font=\Large\bfseries, text=eleveazul] at (5.0,4.35)
    {deslizamento para a esquerda};
  \draw[fisica superficie] (0.65,1.25) -- (9.45,1.25);
  \node[fisica corpo, minimum width=2.0cm, minimum height=1.35cm] (B)
    at (5.0,1.95) {};
  \draw[fisica movimento] ($(B.north)+(0,0.55)$) -- ++(-2.65,0)
    node[fisica rotulo, left] {$\vec v$};
  \draw[fisica forca] (B.center) -- ++(2.65,0)
    node[fisica rotulo, right] {$\vec F_{at}$};
\end{tikzpicture}

\end{document}
